\documentclass[12pt]{article}
\usepackage[utf8]{inputenc}
\usepackage[margin=1in]{geometry}
\usepackage{amsmath}
\usepackage{amsfonts}
\usepackage{amssymb}
\usepackage{enumitem}

\begin{document}

\noindent
CSE400: Fundamentals of Probability in Computing \\
Lecture 20: Linear Transformations and Expectations of Multiple RVs \\
Instructor: Dhaval Patel, PhD \\
Date: March 19, 2025 \\
Institution: Ahmedabad University 

\section*{1. Lecture Outline}
Statistical Representation of Random Vectors \\
Mean Vector, Correlation, and Covariance Matrices \\
Linear Transformations \\
Mapping properties under $\underline{Y}=\underline{A}\underline{X}+\underline{b}$ \\
Multivariate Gaussian Distributions \\
From 1D/2D cases to N-dimensional Joint PDFs \\
Special Cases and Properties \\
Implications of zero correlation and matrix factorization 

\section*{2. Statistical Representation of Random Vectors}
\subsection*{2.1 Random Vector Definition}
For random variables $X_{1}, X_{2}, \dots, X_{N}$, the representation as a vector is called a random vector, denoted as $\underline{X}$. \\
It provides a compact representation of multiple random variables. \\
Vector form: $\underline{X} = \begin{bmatrix} X_{1} \\ X_{2} \\ \vdots \\ X_{N} \end{bmatrix}$.

\subsection*{2.2 Mean Vector}
The mean vector of $\underline{X}$ is defined as $\underline{\mu}_{X}=E[\underline{X}]$. \\
It represents the average value of the random vector. \\
Each element corresponds to the mean of one random variable. \\
Example structure: $\underline{\mu}_{X} = \begin{bmatrix} E[X_{1}] \\ E[X_{2}] \\ \vdots \\ E[X_{N}] \end{bmatrix}$.

\subsection*{2.3 Correlation Matrix}
The correlation matrix represents the correlation between multiple random variables, defined as $R_{XX}=E[\underline{X}\underline{X}^{T}]$. \\
Matrix elements are defined as $R_{XX}(i,j)=E[X_{i}X_{j}]$, where $i$ and $j \le N$. \\
$R_{XX}$ is a symmetric matrix. \\
Full matrix representation: $R_{XX} = \begin{bmatrix} E[X_{1}X_{1}] & E[X_{1}X_{2}] & \dots & E[X_{1}X_{N}] \\ E[X_{2}X_{1}] & \dots & \dots & E[X_{2}X_{N}] \\ \vdots & \vdots & \ddots & \vdots \\ E[X_{N}X_{1}] & E[X_{N}X_{2}] & \dots & E[X_{N}X_{N}] \end{bmatrix}_{N \times N}$.

\subsection*{2.4 Covariance Matrix}
The covariance matrix acts as a relationship map of random vectors. \\
Mathematical definition for a random vector: $C_{XX}=E[(\underline{X}-\underline{\mu}_{X})(\underline{X}-\underline{\mu}_{X})^{T}]$. \\
Practical structure: $C_{XX} = \begin{bmatrix} Var(X_{1}) & Cov(X_{1},X_{2}) & \dots \\ Cov(X_{2},X_{1}) & Var(X_{2}) & \dots \\ \vdots & \vdots & \ddots \end{bmatrix}$. \\
Symmetry property: $Cov(X_{i},X_{j})=Cov(X_{j},X_{i})$, which makes the matrix a "mirror" of itself.

\section*{3. Linear Transformations}
\subsection*{3.1 The Linear Model}
Consider the transformation: $\underline{Y}=\underline{A}\underline{X}+\underline{b}$. \\
$\underline{A}$ represents the transformation matrix, which scales and rotates the data. \\
$\underline{b}$ represents the constant vector, which shifts or translates the center. \\
$\underline{X}$ represents the input random vector. \\
For N-dimensional vectors, the system of equations is: \\
$Y_{1}=a_{11}X_{1}+a_{12}X_{2}+\dots+a_{1N}X_{N}+b_{1}$ \\
$Y_{N}=a_{N1}X_{1}+a_{N2}X_{2}+\dots+a_{NN}X_{N}+b_{N}$ 

\subsection*{3.2 Summary of Properties under Transformation}
If $\underline{Y}=\underline{A}\underline{X}+\underline{b}$, the resulting statistical properties are summarized below:

\vspace{0.5cm}
\begin{tabular}{|l|l|}
\hline
\textbf{Property} & \textbf{Transformed Result} \\
\hline
Mean & $\mu_{Y}=\underline{A}\mu_{X}+\underline{b}$ \\
\hline
Covariance & $C_{YY}=\underline{A}C_{XX}\underline{A}^{T}$ \\
\hline
Correlation & $R_{YY}=\underline{A}R_{XX}\underline{A}^{T}+\underline{A}\mu_{X}\underline{b}^{T}+\underline{b}\mu_{X}^{T}\underline{A}^{T}+\underline{b}\underline{b}^{T}$ \\
\hline
\end{tabular}
\vspace{0.5cm}

\section*{4. Derivations}
\subsection*{4.1 Mean Transformation Derivation}
Definition: $\mu_{Y}=E[\underline{Y}]$. \\
Substitute the linear model: $\mu_{Y}=E[\underline{A}\underline{X}+\underline{b}]$. \\
Apply linearity of expectation: $E[\underline{A}\underline{X}]+E[\underline{b}]$. \\
Final Result: $\mu_{Y}=\underline{A}\mu_{X}+\underline{b}$.

\subsection*{4.2 Correlation Matrix Transformation Derivation}
Definition: $R_{YY}=E[\underline{Y}\underline{Y}^{T}] = E[(\underline{A}\underline{X}+\underline{b})(\underline{A}\underline{X}+\underline{b})^{T}]$. \\
Expansion of terms: $R_{YY}=E[\underline{A}\underline{X}\underline{X}^{T}\underline{A}^{T}+\underline{A}\underline{X}\underline{b}^{T}+\underline{b}\underline{X}^{T}\underline{A}^{T}+\underline{b}\underline{b}^{T}]$. \\
Applying Expectation yields the final result: $R_{YY}=\underline{A}R_{XX}\underline{A}^{T}+\underline{A}\mu_{X}\underline{b}^{T}+\underline{b}\mu_{X}^{T}\underline{A}^{T}+\underline{b}\underline{b}^{T}$.

\subsection*{4.3 Covariance Matrix Transformation Derivation}
Key Observation: The deviation from the mean is independent of the shift $\underline{b}$. \\
$\underline{Y}-\mu_{Y}=(\underline{A}\underline{X}+\underline{b})-(\underline{A}\mu_{X}+\underline{b})=\underline{A}(\underline{X}-\mu_{X})$. \\
Derivation steps: \\
$C_{YY}=E[(\underline{Y}-\mu_{Y})(\underline{Y}-\mu_{Y})^{T}]$. \\
$C_{YY}=E[\underline{A}(\underline{X}-\mu_{X})(\underline{X}-\mu_{X})^{T}\underline{A}^{T}]$. \\
Simplified Result: $C_{YY}=\underline{A}C_{XX}\underline{A}^{T}$. \\
Notice that the shift $\underline{b}$ does not affect the covariance, meaning it does not affect the spread.

\section*{5. Gaussian Distributions}
\subsection*{5.1 1D and 2D Gaussian Distributions}
Univariate Gaussian (1D): For a single random variable $X$, the probability density function (PDF) is $f_{X}(x)=\frac{1}{\sqrt{2\pi\sigma^{2}}}exp(-\frac{(x-\mu)^{2}}{2\sigma^{2}})$. \\
Joint Gaussian (2D): For two random variables $X$ and $Y$, the joint PDF is $f_{XY}(x,y)=\frac{exp\{-\frac{1}{2(1-\rho_{XY}^{2})}[(\frac{x-\mu_{X}}{\sigma_{X}})^{2}-2\rho_{XY}\frac{(x-\mu_{X})(y-\mu_{Y})}{\sigma_{X}\sigma_{Y}}+(\frac{y-\mu_{Y}}{\sigma_{Y}})^{2}]\}}{2\pi\sigma_{X}\sigma_{Y}\sqrt{1-\rho_{XY}^{2}}}$.

\subsection*{5.2 Multivariate Gaussian Distribution}
For an N-dimensional random vector $\underline{X}=[X_{1}, X_{2}, \dots, X_{N}]^{T}$. \\
The General Vector PDF is defined as: $f_{\underline{X}}(x)=\frac{1}{\sqrt{(2\pi)^{N}det(C_{XX})}}exp(-\frac{1}{2}(x-\underline{\mu}_{X})^{T}C_{XX}^{-1}(x-\underline{\mu}_{X}))$. \\
Structure Components: $\underline{\mu}_{X}$ is the Mean vector of size $(N \times 1)$, and $C_{XX}$ is the Covariance matrix of size $(N \times N)$.

\subsection*{5.3 Special Case: Uncorrelated Gaussian Random Variables}
For mutually uncorrelated Gaussian random variables, $Cov(X_{i},X_{j})=0$ for all $i \ne j$. \\
The covariance matrix simplifies to a diagonal matrix, which implies the inverse $C_{XX}^{-1}$ is also a diagonal matrix containing elements $\sigma_{n}^{-2}$ on the diagonal. \\
Substituting the inverse back into the density function yields the exponent: $exp(-\frac{1}{2}\sum_{n=1}^{N}(\frac{x_{n}-\mu_{n}}{\sigma_{n}})^{2})$. \\
Interpretation: The joint Gaussian PDF factorizes into the product of individual Gaussian PDFs: $f_{\underline{x}}(\underline{x})=\prod_{n=1}^{N}\frac{1}{\sqrt{2\pi\sigma_{n}^{2}}}exp(-\frac{(x_{n}-\mu_{n})^{2}}{2\sigma_{n}^{2}})$.

\sectionsection*{6. Worked Examples}
\subsection*{6.1 Numerical Linear Transformation}
Given the linear transformation $\underline{Y}=\underline{A}\underline{X}$ where: \\
$\underline{\mu}_{X} = \begin{bmatrix} 2 \\ -1 \\ 1 \end{bmatrix}$, $\underline{A} = \begin{bmatrix} 1 & 0 & -1 \\ 0 & -1 & 1 \\ -1 & 1 & 0 \end{bmatrix}$, and $R_{XX} = \begin{bmatrix} 13 & 2 & 3 \\ 2 & 10 & 3 \\ 3 & 3 & 10 \end{bmatrix}$.

\noindent
Find Mean Vector of Y: \\
$\underline{\mu}_{Y} = \underline{A}\underline{\mu}_{X} = \begin{bmatrix} 1 & 0 & -1 \\ 0 & -1 & 1 \\ -1 & 1 & 0 \end{bmatrix} \begin{bmatrix} 2 \\ -1 \\ 1 \end{bmatrix} = \begin{bmatrix} 1 \\ 2 \\ -3 \end{bmatrix}$.

\noindent
Find Correlation Matrix of Y: \\
Step 1: Compute $\underline{A}R_{XX} = \begin{bmatrix} 1 & 0 & -1 \\ 0 & -1 & 1 \\ -1 & 1 & 0 \end{bmatrix} \begin{bmatrix} 13 & 2 & 3 \\ 2 & 10 & 3 \\ 3 & 3 & 10 \end{bmatrix} = \begin{bmatrix} 10 & -1 & -7 \\ 1 & -7 & 7 \\ -11 & 8 & 0 \end{bmatrix}$. \\
Step 2: Multiply by $\underline{A}^{T}$ to get $R_{YY} = \begin{bmatrix} 10 & -1 & -7 \\ 1 & -7 & 7 \\ -11 & 8 & 0 \end{bmatrix} \begin{bmatrix} 1 & 0 & -1 \\ 0 & -1 & 1 \\ -1 & 1 & 0 \end{bmatrix} = \begin{bmatrix} 17 & -6 & -11 \\ -6 & 14 & -8 \\ -11 & -8 & 19 \end{bmatrix}$.

\noindent
Find Covariance Matrix of Y: \\
Use property $C_{YY} = R_{YY} - \underline{\mu}_{Y}\underline{\mu}_{Y}^{T}$. \\
Outer product of the mean: $\underline{\mu}_{Y}\underline{\mu}_{Y}^{T} = \begin{bmatrix} 1 \\ 2 \\ -3 \end{bmatrix} \begin{bmatrix} 1 & 2 & -3 \end{bmatrix} = \begin{bmatrix} 1 & 2 & -3 \\ 2 & 4 & -6 \\ -3 & -6 & 9 \end{bmatrix}$. \\
Subtraction: $C_{YY} = \begin{bmatrix} 17 & -6 & -11 \\ -6 & 14 & -8 \\ -11 & -8 & 19 \end{bmatrix} - \begin{bmatrix} 1 & 2 & -3 \\ 2 & 4 & -6 \\ -3 & -6 & 9 \end{bmatrix} = \begin{bmatrix} 16 & -8 & -8 \\ -8 & 10 & -2 \\ -8 & -2 & 10 \end{bmatrix}$. (Note: Source 283 shows an intermediate matrix with -9 in corners, corrected to -11 based on $R_{YY}$).

\subsection*{6.2 2D Joint Gaussian PDF Derivation}
Given mean vector $\underline{\mu}_{X} = \begin{bmatrix} \mu_{1} \\ \mu_{2} \end{bmatrix}$ and covariance matrix $C_{XX} = \begin{bmatrix} \sigma_{1}^{2} & \rho\sigma_{1}\sigma_{2} \\ \rho\sigma_{1}\sigma_{2} & \sigma_{2}^{2} \end{bmatrix}$. \\
Determinant computation: $det(C_{XX}) = \sigma_{1}^{2}\sigma_{2}^{2} - (\rho\sigma_{1}\sigma_{2})^{2} = \sigma_{1}^{2}\sigma_{2}^{2}(1-\rho^{2})$. \\
Inverse of covariance matrix: $C_{XX}^{-1} = \frac{1}{det(C_{XX})}adj(C_{XX}) = \frac{1}{1-\rho^{2}}\begin{bmatrix} \sigma_{1}^{-2} & -\rho\sigma_{1}^{-1}\sigma_{2}^{-1} \\ -\rho\sigma_{1}^{-1}\sigma_{2}^{-1} & \sigma_{2}^{-2} \end{bmatrix}$. \\
Evaluate quadratic form $Q(x) = (\underline{x}-\underline{\mu}_{X})^{T}C_{XX}^{-1}(\underline{x}-\underline{\mu}_{X})$. \\
After algebraic simplification, this yields: $Q(x) = \frac{1}{1-\rho^{2}}[(\frac{x_{1}-\mu_{1}}{\sigma_{1}})^{2} - 2\rho(\frac{x_{1}-\mu_{1}}{\sigma_{1}})(\frac{x_{2}-\mu_{2}}{\sigma_{2}}) + (\frac{x_{2}-\mu_{2}}{\sigma_{2}})^{2}]$.

\section*{7. Case Study: Smart City Traffic + Weather System (Ahmedabad)}
\subsection*{7.1 System Definition}
Design of an intelligent system predicting traffic congestion using transportation and climate variables. \\
Define a random vector $\underline{X}$ of size $N=5$: \\
$X_{1}$: Traffic flow (vehicles/min), mean = 120. \\
$X_{2}$: Average vehicle speed (km/h), mean = 35. \\
$X_{3}$: Rainfall intensity (mm/hr), mean = 2. \\
$X_{4}$: Temperature ($^\circ$C), mean = 32. \\
$X_{5}$: Air pollution index (AQI), mean = 180.

\subsection*{7.2 Dependencies (Covariance Map)}
The covariance matrix acts as the dependency map of the city. By definition and intuition: \\
$Cov(X_{1}, X_{3}) > 0$: More rain implies more congestion (traffic increases). \\
$Cov(X_{2}, X_{3}) < 0$: More rain implies lower speed. \\
$Cov(X_{1}, X_{5}) > 0$: More traffic implies higher pollution. \\
$Cov(X_{4}, X_{5}) < 0$: Higher temperature may cause pollution to disperse.

\subsection*{7.3 Probabilistic Modeling}
Assume the entire system follows a Joint Gaussian Distribution: $\underline{X} \sim \mathcal{N}(\mu, \Sigma)$. \\
Gaussian assumption is used due to the Central Limit Theorem (aggregation of many small effects) and because sensor noise is often Gaussian. \\
This model allows for easier analytics via closed-form solutions.

\subsection*{7.4 Transformation into Congestion Index}
Suppose the city defines a derived metric, the Congestion Index, via linear transformation: $Y = 0.6X_{1} - 0.4X_{2} + 0.8X_{3}$. \\
Since the input is Gaussian, the linear transformation $Y$ is also Gaussian. \\
The mean of $Y$ is given by $E[Y]=a^{T}\mu$, and the variance is $Var(Y)=a^{T}\Sigma a$.

\section*{8. Homework Assignments}
\subsection*{8.1 Finding the Mean Vector}
Problem Statement: Given a random vector $\underline{X}$ with matrices $C_{XX} = \begin{bmatrix} 5 & -1 & 2 \\ -1 & 5 & -2 \\ 2 & -2 & 8 \end{bmatrix}$ and $R_{XX} = \begin{bmatrix} 6 & 1 & 1 \\ 1 & 9 & -4 \\ 1 & -4 & 9 \end{bmatrix}$, find the mean vector $\underline{\mu}_{X}$. 

\noindent
Hint: Use the relationship $C_{XX}=R_{XX}-\underline{\mu}_{X}\underline{\mu}_{X}^{T}$ and rearrange to solve for the outer product: $\underline{\mu}_{X}\underline{\mu}_{X}^{T} = R_{XX} - C_{XX}$. Identify elements of $\underline{\mu}_{X}$ from the diagonal of the result.

\subsection*{8.2 Variance Proof}
Homework: Prove the following variance property: $Var[X+Y] = Var[X] + Var[Y] + 2Cov(X,Y)$.

\end{document}